% !TeX root = ../../Book5.tex
% 纲要标题：### 21.2 形式函数定理的连续证明
% 纲要位置：几何学计划纲要.md，第 5864 行。
% 正文状态：未撰写；本文件目前是框架。

% 纲要内容开始
% 
% * 比较映射及主版本
% 
%   对 Noether 环 A、理想 I、固有 f:X→Spec(A)、凝聚 F，构造 H^q(X,F)^∧→lim_n H^q(X,F/I^nF)。
% 
% * 诱导滤过的等价
% 
%   证明由上同调映射核得到的滤过与 I 进滤过互相控制。
% 
% * 核的稳定与有限性
% 
%   用有限分次模及 Artin–Rees 型估计控制截断产生的核。
% 
% * 像系统的 Mittag–Leffler 性
% 
%   不把逆极限当作一般正合函子，单独证明所用系统的稳定像。
% 
% * 极限上的同构
% 
%   分别验证单射与满射，并识别两侧的拓扑。
% 
% * 局部化与纤维邻域
% 
%   得到高阶直像茎完备化和无穷小纤维上同调的比较。
% 
% 纲要内容结束

\section{形式函数定理的连续证明}
\label{g5:sec:2102}

待撰写

\subsection{比较映射及主版本}
\label{g5:subsec:2102:01}
待撰写

\subsection{诱导滤过的等价}
\label{g5:subsec:2102:02}
待撰写

\subsection{核的稳定与有限性}
\label{g5:subsec:2102:03}
待撰写

\subsection{像系统的 Mittag–Leffler 性}
\label{g5:subsec:2102:04}
待撰写

\subsection{极限上的同构}
\label{g5:subsec:2102:05}
待撰写

\subsection{局部化与纤维邻域}
\label{g5:subsec:2102:06}
待撰写
% 纲要条目排为编号小节标题，并进入目录与书签；说明段仍保留于纲要注释中。
